\setcounter{definition}{0}
\setcounter{theorem}{0}
\subsection{Bevezetés}
\subsubsection{Fogalma}
\begin{remark*}
	Jelölés: $\mathbb{K} \in \{ \mathbb{C, R, Q, Z, Z}_p \}$
\end{remark*}
\begin{definition}
	A $\mathbb{K}$ feletti polinomok halmaza az $x$ és $\mathbb{K}$ elemei által
	az $+, -, \cdot$ segítségével alkotott formális kifejezések.
	\begin{equation*}
		\mathbb{K}[x] = \{
			c_n \cdot x^n + \dots + c_0: n \geq 0, c_n, \dots c_0 \in \mathbb{K}
		\}
	\end{equation*}
\end{definition}
\begin{remark*}
	Jelölése: $\mathbb{K}[x]$
\end{remark*}
\begin{definition}
	Adott $f = c_n \cdot x^n + \dots + c_0$ polinom
	együtthatói a $c_n, \dots, c_0$ számok, míg $c_n \neq 0$ esetén
	$f$ foka $\deg f = n$ és főegyütthatója $c_n$.
\end{definition}
